\section*{Chapitre \ref{ch:b1:predominant-reaction} --- Équilibres acido-basiques et méthode de la réaction prépondérante}

\begin{solution}{exo:b1:predominant-reaction:1}
Bases conjuguées~: \ce{HSO4-}, \ce{CO3^{2-}}, \ce{NH3}, \ce{OH-}. Acides
conjugués~: \ce{NH4+}, \ce{H2PO4-}, \ce{H2O}, \ce{H3O+}. \omterm{def:b1:predominant-reaction:ampholyte}{Ampholytes}~:
\ce{HCO3-}, \ce{H2O}, \ce{HPO4^{2-}} (et \ce{HSO4-}, base de l'acide
sulfurique et acide de son second couple).
\end{solution}

\begin{solution}{exo:b1:predominant-reaction:2}
\ce{HF} ($K_a = \num{6.9e-4}$) $>$ \ce{CH3COOH} (\num{1.7e-5}) $>$ \ce{HClO}
(\num{2.8e-8}) $>$ \ce{NH4+} (\num{5.6e-10}).
\end{solution}

\begin{solution}{exo:b1:predominant-reaction:3}
\omterm{def:b1:predominant-reaction:strong-acid}{Acide fort}~: $\mathrm{pH} = 2.00$. \omterm{def:b1:predominant-reaction:strong-acid}{Base forte}~: $[\ce{OH-}] = 10^{-2}$,
$\mathrm{pH} = 14.00 - 2.00 = 12.00$.
\end{solution}

\begin{solution}{exo:b1:predominant-reaction:4}
\ce{CH3COOH} en dessous de pH 4.76, \ce{CH3COO-} au-dessus. À pH 3, la
forme acide prédomine (98\,\%)~; dans le sang (pH 7.4), c'est l'ion
éthanoate. À pH 3.76, la fraction acide vaut $1/(1 + 10^{-1}) = 0.91$.
\end{solution}

\begin{solution}{exo:b1:predominant-reaction:5}
\omterm{def:b1:predominant-reaction:predominant-reaction}{Réaction prépondérante} \ce{NH3 + H2O <=> NH4+ + OH-}, $K = 10^{-(14.00 -
9.25)} = 10^{-4.75}$. Si peu de \ce{NH3} réagit~: $[\ce{OH-}] = \sqrt{K C} =
\qty{1.3e-3}{mol/L}$, $\mathrm{pH} = 11.13$. Vérification~: $[\ce{NH4+}] = 1.3\,\%$
de $C$, et
$\mathrm{pH} > \mathrm{p}K_a + 1$~: hypothèse validée.
\end{solution}

\begin{solution}{exo:b1:predominant-reaction:6}
$\mathrm{pH} = \frac12(4.76 + 2.00) = 3.38$ (la valeur exacte est
3.39)~; $\alpha = 10^{-3.38}/0.010 = 0.042$~: 4\,\% de l'acide est
dissocié, moins de 10\,\%, et $\mathrm{pH} < \mathrm{p}K_a - 1$~:
hypothèse validée.
\end{solution}

\begin{solution}{exo:b1:predominant-reaction:7}
\ce{HF + OH- -> F- + H2O}, $K = 10^{14.00 - 3.16} = 10^{10.84}$~:
réaction totale. La \omterm{def:b1:predominant-reaction:predominant-reaction}{solution équivalente} contient \qty{0.050}{mol/L} de
\ce{HF} et autant de \ce{F-}, un tampon~: $\mathrm{pH} = \mathrm{p}K_a = 3.16$.
\end{solution}

\begin{solution}{exo:b1:predominant-reaction:8}
$\mathrm{pH} = 4.76$. L'acide chlorhydrique transforme \qty{0.010}{mol}
d'éthanoate en acide~: $\mathrm{pH} = 4.76 + \log(0.040/0.060) = 4.58$,
soit une variation de 0.18. Dans l'eau pure, le pH passerait de 7.00 à
2.00.
\end{solution}

\begin{solution}{exo:b1:predominant-reaction:9}
Dans l'eau, les deux acides réagissent totalement pour donner
\ce{H3O+}~: le solvant les nivelle. L'acide éthanoïque est une base
beaucoup plus faible que l'eau~; les deux acides n'y réagissent donc
que partiellement, dans des proportions différentes, et l'on peut
comparer leurs forces. L'acide le plus fort qui puisse exister dans
l'eau est \ce{H3O+}.
\end{solution}

\begin{solution}{exo:b1:predominant-reaction:10}
\ce{CO3^{2-} + H2O <=> HCO3- + OH-}, $K = 10^{-(14.00 - 10.33)} = 10^{-3.67}$;
$[\ce{OH-}] = \sqrt{10^{-3.67} \times 0.10} = \qty{4.6e-3}{mol/L}$
(4.6\,\% de $C$), $\mathrm{pH} = 11.67$. La seconde basicité,
\ce{HCO3- + H2O <=>
CO2 + H2O + OH-}, a pour constante $K = 10^{-7.63}$~: elle donne
$[\ce{CO2}] \approx
10^{-7.6}$, négligeable.
\end{solution}

\begin{solution}{exo:b1:predominant-reaction:11}
$[\mathrm{A}] = [\ce{H3O+}] = \alpha C$ et $[\mathrm{HA}] = (1 - \alpha)C$,
donc $K_a = \alpha^2C/(1 - \alpha)$. Avec $C = 10^{-3}$~: $\alpha^2 +
0.0174\,\alpha - 0.0174 = 0$, $\alpha = 0.12$. L'acide est dissocié à
12\,\%~: l'hypothèse d'un acide peu dissocié ($\alpha <
10\,\%$) est en défaut, et il faut résoudre l'équation du second degré.
\end{solution}

\begin{solution}{exo:b1:predominant-reaction:12}
L'acide chlorhydrique impose $[\ce{H3O+}] = \qty{0.010}{mol/L}$~: pH 2.00. Alors
\[
  [\ce{CH3COO-}] = \frac{K_a[\ce{CH3COOH}]}{[\ce{H3O+}]} = \frac{1.74 \times 10^{-5}
  \times 0.10}{0.010} = \qty{1.7e-4}{mol/L},
\]
négligeable devant 0.010~: les ions
\ce{H3O+} apportés par l'\omterm{def:b1:predominant-reaction:strong-acid}{acide fort} repoussent l'équilibre de l'\omterm{def:b1:predominant-reaction:strong-acid}{acide
faible}.
\end{solution}

\begin{solution}{pb:b1:predominant-reaction:1}
\textbf{1.} \ce{H3PO4}/\ce{H2PO4-}, \ce{H2PO4-}/\ce{HPO4^{2-}},
\ce{HPO4^{2-}}/\ce{PO4^{3-}}~; par exemple \ce{H3PO4 <=> H2PO4- + H+}.
\textbf{2.} Frontières à 2.15, 7.21 et 12.34.
\textbf{3.} pH 1~: \ce{H3PO4}~; 5~: \ce{H2PO4-}~; 9~: \ce{HPO4^{2-}}~; 13~:
\ce{PO4^{3-}}.
\textbf{4.} \ce{H2PO4-} et \ce{HPO4^{2-}}.
\textbf{5.} Chaque proton laisse un ion de charge négative plus élevée,
qui retient plus fortement le proton suivant.
\textbf{6.} \ce{H3PO4 + H2O <=> H2PO4- + H3O+}~; la formule
$\frac12(\mathrm{p}K_a + \mathrm{p}C) = 1.58$ viole $\mathrm{pH} <
\mathrm{p}K_a - 1 = 1.15$~: environ un cinquième de l'acide réagit.
\textbf{7.} $x^2 + K_{a1}x - K_{a1}C = 0$ avec $K_{a1} = 10^{-2.15}$~: $x =
\qty{0.0233}{mol/L}$, $\mathrm{pH} = 1.63$.
\textbf{8.} \ce{2H2PO4- <=> H3PO4 + HPO4^{2-}}, $K = 10^{2.15 - 7.21} =
10^{-5.06}$.
\textbf{9.} $\mathrm{pH} = \frac12(2.15 + 7.21) = 4.68$.
\textbf{10.} $\mathrm{pH} = \frac12(7.21 + 12.34) = 9.78$.
\textbf{11.} Aucune~: 1.63, 4.68 et 9.78~; il faut un mélange de deux
d'entre elles.
\textbf{12.} \ce{H3PO4 + OH- -> H2PO4- + H2O}, $K = 10^{14.00 - 2.15} =
10^{11.85}$~: \omterm{def:b1:extent-q-and-k:quantitative}{quantitative}.
\textbf{13.} \qty{0.100}{mol/L} de \ce{H2PO4-}~: pH 4.68.
\textbf{14.} Les \qty{0.050}{mol} d'hydroxyde suivantes en transforment
la moitié~: \qty{0.050}{mol/L} de \ce{H2PO4-} et de \ce{HPO4^{2-}}, pH 7.21.
\textbf{15.} \qty{0.100}{mol/L} de \ce{HPO4^{2-}}~: pH 9.78.
\textbf{16.} Le pH part de 1.63, fait un saut autour de $n = 0.100$
(jusqu'à 4.68), monte lentement en passant par 7.21 à $n = 0.150$, fait
un saut autour de $n = 0.200$ (jusqu'à 9.78), passe 12.34 vers
$n = 0.250$ et atteint environ 12.6 à $n = 0.300$ (\qty{0.100}{mol/L}
de \ce{PO4^3-}~: $x^2/(0.100 - x) = 10^{-1.66}$, $x = 0.037$, pH 12.57).
\textbf{17.} Elle contient des quantités égales de l'acide et de la base
d'un couple de $\mathrm{p}K_a$ 7.21 (\cref{prop:b1:predominant-reaction:buffer}).
\textbf{18.} $10^{7.20 - 7.21} = 0.977$.
\textbf{19.} $n(\ce{H2PO4-}) = 0.100 - x$, $n(\ce{HPO4^{2-}}) = x$.
\textbf{20.} $x/(0.100 - x) = 0.977$, $x = \qty{0.0494}{mol}$.
\textbf{21.} L'acide transforme \qty{0.0010}{mol} de \ce{HPO4^{2-}}~:
$\mathrm{pH} = 7.21 + \log(0.0484/0.0516) = 7.18$, soit une variation de
$-0.02$. Dans l'eau pure~: de 7.00 à 3.00.
\textbf{22.} Non~: le pH dépend du rapport des quantités, que la
dilution ne change pas (tant que les \omterm{def:b1:extent-q-and-k:activity}{activités} restent proches des
concentrations).
\textbf{23.} $0.100 + 0.0494 = \qty{0.149}{mol}$ d'hydroxyde~:
\textbf{\qty{149}{mL}} de la solution à \qty{1.00}{mol/L}, complétés à
\qty{1.00}{L}.
\end{solution}
